\section{Was das Plasma verbessert oder verschlechtert: Auswertung des Datensatzes}\label{sec:qi_einfluss}

\textbf{Lesehilfe.} Die QI-Güte (Quasi-Isodynamie) misst, wie gut ein Plasma die Teilchen im Magnetfeld hält; ein
kleinerer Wert ist besser. Die \emph{Rangkorrelation} sagt nur, ob zwei Größen im Datensatz gemeinsam steigen
(positiv) oder gegenläufig sind (negativ). Sie ist \emph{keine Ursache}: Wir haben nie ein Plasma verändert und die Wirkung gemessen.

Der Nutzer fragte, was sich positiv und was negativ auf das Plasma auswirkt. Die Plasmaform stammt aus dem öffentlichen
Proxima-Datensatz; unsere Kette bewertet nur Spulen und Tragwerk dazu.

\subsection*{Erwartung vorher}
Aus den 182 vorsortierten Spulensätzen (Kraft-Scan) war das Aspektverhältnis der stärkste Einfluss (+0,30, größer ist
schlechter). Erwartet wurde dasselbe Vorzeichen im ganzen Datensatz, mit größerem Betrag (0,3 bis 0,5).

\subsection*{Ergebnis}
158\,685 von 182\,222 Plasmen haben ein gültiges QI. Rangkorrelation mit $\log_{10}$(QI):

\begin{tabular}{lr}
\hline
Merkmal & Rangkorrelation \\ \hline
Randwert der Rotationstransformation (je Periode) & $+0{,}35$ \\
kleinste normierte Gradientenlänge & $-0{,}34$ \\
Aspektverhältnis & $-0{,}31$ \\
größte Elongation & $-0{,}30$ \\
Vakuum-Well & $-0{,}30$ \\
Spiegelverhältnis am Rand & $+0{,}27$ \\
mittlere Dreieckigkeit & $+0{,}22$ \\
Spiegelverhältnis auf der Achse & $-0{,}21$ \\
Flusskompression in schlechter Krümmung & $-0{,}20$ \\
Feldperioden & $+0{,}14$ \\ \hline
\end{tabular}

\smallskip
\noindent Medianes QI nach Feldperioden: 1: 0,0109; 2: 0,0085; 3: 0,0204; 4: 0,0145; 5: 0,0179. Das beste Einzelplasma hat
2 Perioden (QI $6\cdot10^{-5}$).

\subsection*{Vergleich mit der Erwartung: verfehlt}
Das Vorzeichen beim Aspektverhältnis \emph{kehrt sich um} ($-0{,}31$ statt $+0{,}30$). Die Nachprüfung in Bereichen des
Aspektverhältnisses zeigt, warum: unter 10 ist ein größeres Aspektverhältnis besser (etwa $-0{,}17$ bis $-0{,}34$), über 10
schlechter ($+0{,}25$ und mehr). Unsere 182 Sätze liegen alle bei 9 bis 11, dort hat die Korrelation das andere
Vorzeichen. Auch bei den Feldperioden ist die Gesamtzahl ($+0{,}14$) irreführend: innerhalb jeder Periodenzahl
sind es $-0{,}14$ bis $-0{,}37$.

\subsection*{Deutung und Grenzen}
\begin{itemize}\itemsep1pt
\item Die Korrelationen sind schwach (höchstens 0,35) und hängen vom betrachteten Bereich ab. Aus ihnen folgt
  \emph{keine} Regel wie „flacher ist besser“.
\item Belastbar ist nur: Die besten Plasmen liegen bei wenigen Feldperioden (2); 3 Perioden sind im Median am schlechtesten.
\item Die frühere Aussage, mehr Kraft gehe mit besserem QI einher ($-0{,}26$), gilt nur für die 182 Sätze und ist im
  Gesamtdatensatz nicht geprüft.
\item QI ist ein Gütemaß der Feldform. Einschluss und Stabilität des Plasmas sind damit nicht gerechnet.
\end{itemize}
Quelle: \texttt{stellarator/qi\_einfluss.py}, Ergebnis \texttt{docs/recherche/qi\_tragwerk/qi\_einfluss.md}.

\subsection*{Nachtrag: die Merkmale beeinflussen sich gegenseitig}
Der Nutzer vermutete, dass es sich um Bereiche handelt, die sich gegenseitig beeinflussen. Ein gemeinsames Modell
(Gradient Boosting, auf zurückgehaltenen Daten geprüft) bestätigt das: Es erklärt die QI-Güte mit $R^2=0{,}90$, jede
Einzelkorrelation höchstens $0{,}12$. Am wichtigsten sind die beiden Spiegelverhältnisse (auf der Achse und am Rand), die
einzeln gegenläufige Vorzeichen hatten und sich dadurch aufhoben. \emph{Lesehilfe:} Das Spiegelverhältnis sagt, wie stark
das Feld entlang einer Feldlinie schwankt; bei guter QI sind die Schwankungen auf allen Flächen ähnlich.

Bereiche der besten 1 \% (QI höchstens $0{,}00084$; 5. bis 95. Perzentil, in Klammern alle Plasmen):
Aspektverhältnis 7,5 bis 9,7 (5,3 bis 10,6); Elongation 5,3 bis 7,4; Spiegelverhältnis Achse 0,28 bis 0,40 (0,12 bis 0,58),
am Rand 0,31 bis 0,43 (0,20 bis 0,85); Rotationstransformation am Rand 0,07 bis 0,24; höchstens 4 Feldperioden.

\emph{Erwartung vorher:} $R^2>0{,}5$ (getroffen), wichtigstes Merkmal Rotationstransformation oder Gradientenlänge
(verfehlt). \emph{Grenze:} Die Spiegelverhältnisse sind selbst Messgrößen der Feldform und eng mit QI verwandt, also
keine frei wählbaren Entwurfshebel; es bleibt eine Korrelation. Quelle: \texttt{stellarator/qi\_bereiche.py}.

\subsection*{Parameterkorridor für ein optimales Plasma}
Als „gut“ gilt hier QI $\le 0{,}001$ (2420 von 158\,685 Plasmen, 1,5\,\%). Ein \emph{Kasten} aus den 5.\ bis 95.\ Perzentilen der guten Plasmen:
Aspektverhältnis 7,6 bis 9,7; Elongation 5,35 bis 7,45; Spiegelverhältnis Achse 0,27 bis 0,40, Rand 0,30 bis 0,43; Rotationstransformation je Periode
am Rand 0,07 bis 0,25; Vakuum-Well $-0{,}145$ bis $-0{,}01$. Im Kasten sind 23\,\% der Plasmen gut (Grundrate 1,5\,\%) und er enthält 46\,\% aller guten.
Ein flacher Entscheidungsbaum zeigt die Wechselwirkung: Der beste Bereich (Elongation $>6{,}06$, Spiegelverhältnis am Rand $\le0{,}43$, auf der Achse $>0{,}325$,
Aspektverhältnis $\le9{,}09$) hat 40\,\% Reinheit bei 31\,\% Abdeckung. \emph{Erwartung vorher:} Kasten 3 bis 15\,\%, Baum 20 bis 50\,\% (getroffen bzw.\ übertroffen).
Unsere Sätze DjLe3eL6, DYU und DPrcxq66 liegen im Kasten, DP9 liegt in mehreren Merkmalen außerhalb.
\emph{Grenzen:} Die Kennzahlen sind Messgrößen der Feldform, keine Entwurfsparameter; auch im Korridor sind drei von fünf Plasmen nicht gut. Quelle: \texttt{stellarator/qi\_korridor.py}.

\subsection*{Von den Kennzahlen zur Form: Formkoeffizienten}
Eingabe sind die Zahl der Feldperioden und die Tabelle der Formkoeffizienten (\texttt{r\_cos}, \texttt{z\_sin}, je $5\times9$ Zahlen); sie beschreiben die Randfläche des Plasmas
als Überlagerung von Wellen. Ein Gradient-Boosting-Modell sagt $\log_{10}$(QI) aus diesen 91 Eingaben mit $R^2=0{,}91$ auf zurückgehaltenen Daten voraus
(Erwartung 0,6 bis 0,85, übertroffen). Am wichtigsten sind die Koeffizienten niedriger Ordnung ($m=0$ und $1$) und die Feldperiodenzahl.
Ein \emph{Kasten} aus Bereichen der fünf wichtigsten Koeffizienten hat nur 2,3\,\% Reinheit (Erwartung 5 bis 15\,\%, verfehlt), der beste Bereich eines Entscheidungsbaums 5\,\%.
\emph{Deutung:} Gute Formen entstehen aus dem Zusammenspiel der Koeffizienten und lassen sich nicht als Bereiche einzelner Koeffizienten angeben; das Modell kann es, ein Kasten nicht.
Der nächste Schritt ist daher, neue Formen im Modell zu suchen und mit dem echten Rechner (VMEC++) nachzuprüfen.
\emph{Grenzen:} nur im Datenbereich gültig, Indexkonvention nicht geprüft, QI misst nur die Feldform. Quelle: \texttt{stellarator/qi\_koeffizienten.py}.

\subsection*{Gegenrechnung mit dem echten Rechner und die offizielle Wertung}
Zehn im Surrogat gefundene Formen wurden mit VMEC++ nachgerechnet (QI braucht die Einstellung \texttt{default\_high\_fidelity}, 70 bis 155\,s je Form). Bei weiter Suche (1,5 Standardabweichungen von den Daten) ließ sich keine der zehn Formen rechnen. Mit kleinen Schritten (höchstens 0,3 Standardabweichungen) rechneten 9 von 10, alle mit QI unter 0,00045, aber nur 3 besser als ihre Startform.
\emph{Erwartung vorher:} 2 bis 4 von 10 unter QI 0,001, höchstens eine bessere als DjLe3eL6; \emph{verfehlt} bei der weiten Suche. Die einzige eigene Form (halber Schritt von der besten Datenform, QI $4{,}7\cdot10^{-5}$ gegen $6{,}2\cdot10^{-5}$) ist nur eine kleine Abwandlung.
Nach der \emph{offiziellen Wertung} der Rangliste (lokal bestimmt, nichts eingereicht) erhalten alle Formen den Wert 0, auch die Datenformen und die eigene: verletzt sind Spiegelverhältnis (0,35 bis 0,54 statt höchstens 0,2 bis 0,25), Rotationstransformation am Rand je Periode (0,07 bis 0,21 statt mindestens 0,25), Vakuum-Well, Flusskompression und Elongation. QI ist dort nur Nebenbedingung.
\emph{Folge:} Eine Form mit niedrigem QI ist nicht automatisch eine gute Form im Sinn der Rangliste. Ein Upload ist daher nicht sinnvoll.
